\section{Classical reductions with one index equal to two}
\label{one2:sec:general}

This section settles the three Gaussian weight-four triple identities
printed in Chapter~4 of the source manuscript.  The proof gives a uniform
formula for every multiple polylogarithm whose indices consist of one
$2$ and any number of $1$'s.  The general phenomenon is classical:
Borwein and Straub give the complementary-argument formula for
$\Li_{2,\{1\}^{p-1}}$ in \cite[Eq.~(33)]{BorweinStraub2015}, including
the weight-four instance in their Eq.~(34).  The results below reorganize
that reduction at the reciprocal complementary argument, make the
dependence on the position of the $2$ explicit, and prove the specific
Gaussian formulas in the manuscript.  We do not claim that classical
reducibility of this family is a new phenomenon.

Throughout this section, a single displayed argument abbreviates the
remaining arguments $1$:
\[
 \Li_{s_1,\ldots,s_d}(z)
 :=\Li_{s_1,\ldots,s_d}(z,1,\ldots,1).
\]
The defining series has strictly decreasing summation indices.  Set
\[
 U_{a,b}(z)=\Li_{\{1\}^{a},2,\{1\}^{b}}(z),\qquad
 H_p(z)=\Li_{2,\{1\}^{p-1}}(z),
\]
where $a,b\ge0$ and $p\ge1$.  The weight of $U_{a,b}$ is
$N=a+b+2$, and its series depth is $N-1$.

\subsection{A logarithmic moment and its branch conventions}

We use principal branches and first work on
\[
 \mathcal D=\mathbb C\setminus[0,\infty),\qquad
 w=\frac1{1-z},\qquad L=\Log w=-\Log(1-z).
\]
On this domain, $w$ avoids the standard polylogarithm cut $[1,\infty)$
and the logarithm cut $(-\infty,0]$.  All terms in the formulas below
are therefore single-valued analytic functions.  The formulas also
extend to the usual slit plane by consistent continuation; no such
extension is needed at $z=i$ or at a nonreal root of unity.

\begin{lemma}[Logarithmic moment]
\label{one2:lem:moment}
Let $q(t)=-\Log(1-t)$.  For an integration path from $0$ to $z$ avoiding
$\{0,1\}$ apart from its initial point $0$, with the branches
continued along that path,
\begin{equation}
 U_{a,b}(z)=\frac1{a!(b+1)!}
 \int_0^z\frac{(q(z)-q(t))^a q(t)^{b+1}}{t}\,dt.
 \label{one2:eq:moment}
\end{equation}
Consequently,
\begin{equation}
 U_{a,b}(z)=\sum_{j=0}^a(-1)^j
 \binom{b+j+1}{j}\frac{q(z)^{a-j}}{(a-j)!}\,H_{b+j+1}(z).
 \label{one2:eq:triangular}
\end{equation}
\end{lemma}
\begin{proof}
The iterated-integral word of $U_{a,b}$ is
$1^a0\,1^{b+1}$, with forms $dt/(1-t)$ and $dt/t$ attached to $1$ and
$0$, respectively.  The innermost $b+1$ identical forms contribute
$q(t)^{b+1}/(b+1)!$.  The $a$ identical forms between $t$ and $z$
contribute $(q(z)-q(t))^a/a!$.  Integrating the single remaining $0$
form proves \eqref{one2:eq:moment}.  Equivalently, one may verify the
same formula by differentiating successively with respect to $z$ and
using the zero initial values at $z=0$.  Expanding the binomial and
using
\[
 H_p(z)=\frac1{p!}\int_0^z\frac{q(t)^p}{t}\,dt
\]
proves \eqref{one2:eq:triangular}.
\end{proof}

\begin{theorem}[Explicit reduction for every position of the $2$]
\label{one2:thm:closed}
For $a,b\ge0$, $N=a+b+2$, and $z\in\mathcal D$,
\begin{align}
 U_{a,b}(z)
 &=\sum_{k=a+1}^{N}(-1)^{a+k}\binom{k-1}{a}
       \frac{L^{N-k}}{(N-k)!}\Li_k(w)
       -\frac{L^N}{N!}\notag\\
 &\quad+(-1)^{b+1}\sum_{j=0}^{a}\binom{b+j+1}{j}
       \frac{L^{a-j}}{(a-j)!}\zeta(b+j+2).
 \label{one2:eq:closed}
\end{align}
In particular,
\begin{equation}
 H_p(z)=(-1)^p\zeta(p+1)-\frac{L^{p+1}}{(p+1)!}
 +\sum_{k=1}^{p+1}(-1)^k
       \frac{L^{p+1-k}}{(p+1-k)!}\Li_k(w).
 \label{one2:eq:height}
\end{equation}
\end{theorem}
\begin{proof}
We first prove \eqref{one2:eq:height} for real $z<0$, so that
$0<w<1$ and every intermediate integral is real.  Substitute
$v=1/(1-t)$ in the logarithmic moment.  Then
\[
 \frac{dt}{t}=\frac{dv}{v(v-1)},\qquad q(t)=\log v,
\]
and
\[
 H_p(z)=\frac1{p!}\int_1^w\frac{\log^p v}{v(v-1)}\,dv.
\]
The differential identity
\begin{equation}
 \frac{d}{dv}\left[
 \sum_{k=1}^{p+1}\frac{(-1)^{k-1}p!}{(p+1-k)!}
       \log^{p+1-k}v\,\Li_k(v)\right]
       =\frac{\log^p v}{1-v}
 \label{one2:eq:primitive}
\end{equation}
follows by cancellation of adjacent terms, using
$\Li_1(v)=-\log(1-v)$ and $\Li'_k(v)=\Li_{k-1}(v)/v$ for $k\ge2$.
Since
$1/[v(v-1)]=-1/v-1/(1-v)$, integration gives
\eqref{one2:eq:height}.  At $v=1$, every term in
\eqref{one2:eq:primitive} with a positive power of $\log v$ tends to
zero.  This includes the apparently singular term
$\log^p v\,\Li_1(v)$, which is
$O((v-1)^p\log(1-v))$ and tends to zero because $p\ge1$.
The sole nonzero endpoint term is $(-1)^pp!\zeta(p+1)$.
Thus the constant in \eqref{one2:eq:height} is fixed, not inferred
from a differential identity alone.  Analytic continuation on the
connected domain $\mathcal D$ proves that identity throughout the
stated domain.

Insert \eqref{one2:eq:height} in
\eqref{one2:eq:triangular}, using $q(z)=L$.  The zeta terms immediately
give the second line of \eqref{one2:eq:closed}.  The coefficient of
$(-1)^k L^{N-k}\Li_k(w)$ is
\[
 \sum_{j=0}^a\frac{(-1)^j\binom{b+j+1}{j}}
                    {(a-j)!(b+j+2-k)!}
 =\frac{(-1)^a\binom{k-1}{a}}{(N-k)!}.
\]
Terms containing a factorial of a negative integer are omitted.  To
verify this finite identity, write its left side as
\[
 \frac{(k-1)!}{(b+1)!a!}
 \sum_{j=0}^a(-1)^j\binom aj\binom{b+j+1}{k-1}
\]
and extract the coefficient of $x^{k-1}$ in
\[
 (1+x)^{b+1}\sum_{j=0}^a(-1)^j\binom aj(1+x)^j
   =(-x)^a(1+x)^{b+1}.
\]
The coefficients vanish when $k\le a$, explaining the lower endpoint
$a+1$ in the sum.  Finally, the coefficient of $-L^N$ equals
\begin{align*}
 \sum_{j=0}^a\frac{(-1)^j\binom{b+j+1}{j}}
                    {(a-j)!(b+j+2)!}
 &=\frac1{a!(b+1)!}\int_0^1t^{b+1}(1-t)^a\,dt\\
 &=\frac1{N!}.
\end{align*}
This proves the complete formula.
\end{proof}

For the Gaussian endpoint the path itself can be inspected explicitly:
\[
 z=ix\quad(0\le x\le1),\qquad
 w(x)=\frac{1+ix}{1+x^2}.
\]
For $x>0$ this path lies in the upper half-plane and does not meet the
polylogarithm cut.  Its initial point is $w(0)=1$, where the endpoint
limits in the proof apply, and its final point is $(1+i)/2$.  This
verifies the branch choices directly, without using an inversion
formula with an unspecified logarithmic correction.

The boundary series at $z=i$ agree with these analytic values.
Indeed the finite nested harmonic sums in a coefficient satisfy a
bound $O((1+\log n)^{d-1})$.  After division by $n^{s_1}$ their
coefficients tend to zero and have summable successive differences.
Summation by parts against the bounded partial sums of $i^n$ proves
convergence, even when $s_1=1$; Abel's theorem then identifies the
series with the radial analytic limit.

\begin{corollary}[A short formula when the $2$ is last]
\label{one2:cor:last}
For every $N\ge2$,
\begin{align}
 U_{N-2,0}(z)
 &=(N-1)\Li_N(w)-L\Li_{N-1}(w)-\frac{L^N}{N!}\notag\\
 &\quad-\sum_{m=0}^{N-2}(N-1-m)\zeta(N-m)\frac{L^m}{m!}.
 \label{one2:eq:last}
\end{align}
Thus only two nonconstant classical polylogarithms are needed in this
subfamily, regardless of series depth.
\end{corollary}
\begin{proof}
Set $a=N-2$, $b=0$ in \eqref{one2:eq:closed} and reindex the zeta sum.
\end{proof}

\begin{corollary}[The top two series-depth layers]
\label{one2:cor:toplayers}
Every single-endpoint multiple polylogarithm of weight $N\ge2$ and
series depth at least $N-1$ has an explicit expression in classical
polylogarithms at $(1-z)^{-1}$, zeta values, and logarithms, on the
domain $\mathcal D$.
\end{corollary}
\begin{proof}
Positive indices with sum $N$ and length $N$ are all $1$, and
$\Li_{\{1\}^N}(z)=L^N/N!$. Those with length $N-1$ have exactly one
index $2$ and all other indices $1$, so Theorem~\ref{one2:thm:closed}
applies. These exhaust the possibilities.
\end{proof}

This statement concerns the literal depth of the defining series.
It does not identify that depth with a motivic filtration or prove a
minimal depth for the resulting numerical constants.

\subsection{One possible new direction at each weight}

Put
\[
 \ell=\log2,\qquad w=\frac{1+i}{2},\qquad
 \mu_k=\Re\Li_k(w),\qquad\lambda_k=\Im\Li_k(w).
\]
For $N\ge2$, let $\mathscr P_N$ be the rational span of products of
at least two positive-weight atoms selected from
\[
 \pi,\ell,\quad \zeta(k),\mu_k,\lambda_k\quad(2\le k<N),
\]
whose assigned weights sum to $N$; the assigned weights are $1$ for
$\pi,\ell$ and $k$ for the atoms with subscript or argument $k$.
This definition is a concrete span of numerical products.  It does
not assume that numerical periods form a direct sum by weight.

\begin{corollary}[Signed binomial coefficients modulo products]
\label{one2:cor:primitive}
For $a+b+2=N$, set
$c_{N,a}=(-1)^{N+a}\binom{N-1}{a}$.  Then
\begin{align}
 U_{a,b}(i)&\equiv c_{N,a}\bigl(\Li_N(w)-\zeta(N)\bigr)
        \pmod{\mathscr P_N+i\mathscr P_N},\label{one2:eq:complexprimitive}\\
 \Re U_{a,b}(i)&\equiv c_{N,a}\bigl(\mu_N-\zeta(N)\bigr)
        \pmod{\mathscr P_N},\label{one2:eq:realprimitive}\\
 \Im U_{a,b}(i)&\equiv c_{N,a}\lambda_N
        \pmod{\mathscr P_N}.\label{one2:eq:imagprimitive}
\end{align}
In particular, the imaginary parts of this entire family contribute
at most one direction after products of the stated lower-weight
atoms have been factored out.
\end{corollary}
\begin{proof}
At $z=i$, $L=-\ell/2+i\pi/4$, and
$\Li_1(w)=-\overline L$.  In \eqref{one2:eq:closed}, every term is a
product of lower-weight atoms except the term $k=N$ and the term
$j=a$ in the zeta sum.  Their respective coefficients are
$c_{N,a}$ and $-c_{N,a}$.  Taking real and imaginary parts proves the
statements.
\end{proof}

The coefficients of the highest-weight half-point imaginary value are
\[
 \begin{array}{c|l}
 N&c_{N,0},\ldots,c_{N,N-2}\\\hline
 4&1,-3,3\\
 5&-1,4,-6,4\\
 6&1,-5,10,-10,5.
 \end{array}
\]
The corollary is an upper bound, and proves neither that $\lambda_N$
is nonzero in the quotient nor that it is independent of constants
from other argument families.  It also does not identify series
depth with an additive motivic depth filtration.

\input{sections/sixth_root}

\section{Proof of the three Gaussian weight-four identities}
\label{one2:sec:triples}

\begin{lemma}[The lower half-point values]
\label{one2:lem:halfpoint}
With $G=\beta(2)$ and the notation above,
\begin{align}
 \Li_2(w)&=\frac{5\pi^2}{96}-\frac{\ell^2}{8}
              +i\left(G-\frac{\pi\ell}{8}\right),
       \label{one2:eq:Li2half}\\
 \mu_3&=\frac{35\zeta(3)}{64}-\frac{5\pi^2\ell}{192}
              +\frac{\ell^3}{48}.
       \label{one2:eq:mu3}
\end{align}
\end{lemma}
\begin{proof}
Differentiation, followed by evaluation at $z=0$, proves Landen's
dilogarithm identity
\[
 \Li_2\left(\frac z{z-1}\right)
       =-\Li_2(z)-\frac12\Log^2(1-z)
\]
along the segment from $0$ to $i$.  At $z=i$ its transformed argument
is $\overline w$.  The defining series gives
$\Li_2(i)=-\pi^2/48+iG$, while
$\Log(1-i)=\ell/2-i\pi/4$.  Taking complex conjugates proves
\eqref{one2:eq:Li2half}.

For completeness, the trilogarithm relation needed here is
\begin{align}
 \Li_3(x)+\Li_3\left(\frac x{x-1}\right)+\Li_3(1-x)
 &=\zeta(3)+\zeta(2)\Log(1-x)\notag\\
 &\quad-\frac12\Log x\,\Log^2(1-x)
       +\frac16\Log^3(1-x).
 \label{one2:eq:trilog}
\end{align}
It holds on the vertical strip $0<\Re x<1$.  Its derivative follows
from the preceding Landen identity and
$\Li_2(x)+\Li_2(1-x)=\zeta(2)-\Log x\Log(1-x)$;
the constant is obtained by letting $x\downarrow0$ through positive
real values.  These dilogarithm identities can themselves be checked
by differentiation and the same endpoint limits, so no
unproved polylogarithm identity is being used.

Now set $x=w$.  Then $1-w=\overline w$, $w/(w-1)=-i$,
$\Log w=L=-\ell/2+i\pi/4$, and
$\Log(1-w)=\overline L$.  Also
$\Re\Li_3(-i)=-3\zeta(3)/32$, directly from its even terms.
Taking real parts of \eqref{one2:eq:trilog} yields
\[
 2\mu_3-\frac{3\zeta(3)}{32}
 =\zeta(3)-\frac{5\pi^2\ell}{96}+\frac{\ell^3}{24},
\]
which is \eqref{one2:eq:mu3}.
\end{proof}

\begin{theorem}[The manuscript's Gaussian triples]
\label{one2:thm:triples}
All three identities below hold exactly:
\begin{align}
 \Im\Li_{2,1,1}(i)
 &=\lambda_4+\frac12\ell\lambda_3
       -\frac{G(\pi^2-4\ell^2)}{32}
       -\frac{\pi(8\ell^3+105\zeta(3))}{768},
       \label{one2:eq:211}\\
 \Im\Li_{1,2,1}(i)
 &=-3\lambda_4-\ell\lambda_3
       +\frac{G(\pi^2-4\ell^2)}{32}
       -\frac{3\pi^3\ell}{256}
       +\frac{\pi(2\ell^3+67\zeta(3))}{128},
       \label{one2:eq:121}\\
 \Im\Li_{1,1,2}(i)
 &=3\lambda_4+\frac12\ell\lambda_3
       +\frac{5\pi^3\ell}{192}
       -\frac{163\pi\zeta(3)}{256}.
       \label{one2:eq:112}
\end{align}
\end{theorem}
\begin{proof}
The case $p=3$ of \eqref{one2:eq:height} is
\begin{equation}
 H_3(i)=\Li_4(w)-L\Li_3(w)+\frac{L^2}{2}\Li_2(w)
        +\frac{L^3\overline L}{6}-\frac{L^4}{24}-\zeta(4).
 \label{one2:eq:H3complex}
\end{equation}
Taking imaginary parts and inserting Lemma~\ref{one2:lem:halfpoint}
gives \eqref{one2:eq:211}; all coefficients are rational arithmetic
in $L=-\ell/2+i\pi/4$.

The case $p=2$ gives the useful explicit intermediary
\begin{equation}
 H_2(i)=\frac{29\zeta(3)-16\pi G}{64}
       +i\left(-\lambda_3-\frac{G\ell}{2}
                    +\frac{\pi\ell^2}{32}+\frac{3\pi^3}{128}\right).
 \label{one2:eq:H2complex}
\end{equation}
The word shuffles, or directly \eqref{one2:eq:triangular}, give
\begin{align*}
 \Li_{1,2,1}(i)&=LH_2(i)-3H_3(i),\\
 \Li_{1,1,2}(i)&=\frac{L^2}{2}\Li_2(i)-2LH_2(i)+3H_3(i).
\end{align*}
Inserting \eqref{one2:eq:H2complex}, \eqref{one2:eq:211}, and
$\Li_2(i)=-\pi^2/48+iG$ proves the last two identities.  This also
explains the coefficients $1,-3,3$ of $\lambda_4$: they are the
weight-four instances of Corollary~\ref{one2:cor:primitive}.
\end{proof}

These formulas can now be moved out of the source manuscript's
numerically supported list and labeled theorems.  Their coefficients
require no correction; the change is a rigorous justification of the
previous proof status.

\subsection{Further exact identities at weights five and six}

The same calculation gives formulas beyond the three displayed
manuscript rows.  They involve the unavoidable output of this
particular reduction, namely higher classical values at $w$; no
independence claim about those values is implied.

\begin{corollary}[Two higher-weight Gaussian reductions]
\label{one2:cor:56}
With $H_4=\Li_{2,1,1,1}$ and $H_5=\Li_{2,1,1,1,1}$,
\begin{align}
 \Im H_4(i)
 &=-\lambda_5-\frac{\ell\lambda_4}{2}+\frac{\pi\mu_4}{4}
       +\frac{\pi^2-4\ell^2}{32}\lambda_3\notag\\
 &\quad+\frac{G\ell(3\pi^2-4\ell^2)}{192}
       +\frac{35\pi\ell\zeta(3)}{512}
       +\frac{\pi\ell^2(\ell^2-\pi^2)}{384}
       -\frac{17\pi^5}{92160},
       \label{one2:eq:H4}\\
 \Im H_5(i)
 &=\lambda_6+\frac{\ell\lambda_5}{2}-\frac{\pi\mu_5}{4}
       +\frac{4\ell^2-\pi^2}{32}\lambda_4
       -\frac{\pi\ell\mu_4}{8}\notag\\
 &\quad+\frac{\ell(4\ell^2-3\pi^2)}{192}\lambda_3
       +\frac{G(16\ell^4-24\ell^2\pi^2+\pi^4)}{6144}\notag\\
 &\quad+\frac{35\pi(\pi^2-12\ell^2)\zeta(3)}{24576}
       +\frac{\pi\ell^3(5\pi^2-3\ell^2)}{5760}.
       \label{one2:eq:H5}
\end{align}
\end{corollary}
\begin{proof}
Use $p=4$ and $p=5$ in \eqref{one2:eq:height}, substitute
$L=-\ell/2+i\pi/4$ and Lemma~\ref{one2:lem:halfpoint}, and take
imaginary parts.  The accompanying coefficient generator performs
these polynomial operations over $\mathbb Q$ and also emits every
position of the $2$ at weights four, five, and six.
\end{proof}

The real parts can be treated by the same formula.  For example,
\begin{equation}
 \Re\Li_{2,1,1}(i)=\mu_4+\frac{\pi\lambda_3}{4}
 +\frac{\pi G\ell}{8}+\frac{35\ell\zeta(3)}{128}
 +\frac{\ell^4}{384}-\frac{11\pi^2\ell^2}{768}
 -\frac{1249\pi^4}{92160}.
 \label{one2:eq:real211}
\end{equation}
The machine-readable tables include both real and imaginary parts,
so this companion family can be reproduced without manual
transcription.

\subsection{Exponential series evaluation and independent checks}

At $w=(1+i)/2$ one has $r=|w|=2^{-1/2}<1$.  Truncating only the
classical series in \eqref{one2:eq:closed} after $M$ terms, while
treating $L$ and the zeta values as exact, produces an error bounded by
\begin{equation}
 \frac{r^{M+1}}{1-r}
 \sum_{k=a+1}^N\binom{k-1}{a}
       \frac{|L|^{N-k}}{(N-k)!(M+1)^k}.
 \label{one2:eq:truncbound}
\end{equation}
Indeed the tail of each $\Li_k(w)$ is bounded by
$r^{M+1}/((M+1)^k(1-r))$, and the triangle inequality gives
\eqref{one2:eq:truncbound}.  Thus the classical series converge
geometrically even when the original depth-$N-1$ boundary series is
only conditionally convergent.  A complete interval implementation
must additionally enclose $L$ and the zeta values; this displayed
bound isolates the polylogarithm truncation error.

The accompanying independent numerical script compares the closed
formula with the logarithmic moment \eqref{one2:eq:moment} for all
positions of the $2$ at weights $2$ through $8$, at the four endpoints
\[
 i,\qquad\frac{-1+i\sqrt3}{2},\qquad 0.3+0.2i,\qquad-0.6.
\]
At $100$ working decimal digits, all $112$ comparisons have residual
below $5.14\cdot10^{-100}$.  Direct nested-series summation at the two
interior points through weight six gives residuals below
$1.46\cdot10^{-100}$.  Separate implementations of the three
manuscript rows and \eqref{one2:eq:H4}--\eqref{one2:eq:H5} agree with
quadrature to below $3.53\cdot10^{-101}$.  These residuals are
numerical cross-checks of formulas already proved above, not
rigorous interval-error certificates or independence evidence.

